\section{The dense FFN bottleneck}
\label{sec:ch01-the-dense-ffn-bottleneck}

Show that the dense feed-forward block is the natural place to introduce sparsity because every token pays for every hidden unit.

\input{figures/ch01/fig-the-dense-ffn-bottleneck}

Use a four-token batch with d\_model=8 and an FFN expansion factor of 4; count multiply-adds for all tokens.

\begin{table}[H]
\centering
\caption{Parameter and activation shape comparison for attention, dense FFN, and the residual path.}
\label{tab:ch01-the-dense-ffn-bottleneck}
\begin{tabularx}{\textwidth}{p{0.22\textwidth}X p{0.22\textwidth}}
\toprule
Item & Planned content & Status \\
\midrule
Mechanism & Show that the dense feed-forward block is the natural place to introduce sparsity because every token pays for every hidden unit. & Placeholder \\
Mini-example & Use a four-token batch with d-model=8 and an FFN expansion factor of 4; count multiply-adds for all tokens. & Placeholder \\
Diagnostic & A small cost table proving that increasing H increases every token path. & Placeholder \\
\bottomrule
\end{tabularx}
\end{table}


The planned mechanics should introduce the equation before the implementation listing.

\begin{equation}
\label{eq:ch01-the-dense-ffn-bottleneck}
\text{Dense FFN cost estimate using D, H, and number of tokens N.}
\end{equation}


\begin{lstlisting}[style=bookpython,caption={Planned code placeholder for the dense FFN bottleneck.},label={lst:ch01-the-dense-ffn-bottleneck}]
def the_dense_ffn_bottleneck_placeholder(x, config):
    """Chapter 1.1 placeholder.

    Planned purpose: Tiny function that estimates dense FFN parameters and per-token matrix multiplies.
    Replace this stub during section development.
    """
    raise NotImplementedError("Implement this section's code path.")
\end{lstlisting}


\begin{checkpointbox}
A small cost table proving that increasing H increases every token path.
\end{checkpointbox}

Once the bottleneck is visible, the next question is whether every token really needs the same FFN.
